\section{Introduction}

Write \(R=\rad(abc)=\prod_{p\mid abc}p\) and \(H=\log c\).
The \(abc\) conjecture predicts \(H\le(1+\varepsilon)\log R+O_\varepsilon(1)\)
for every coprime positive solution of \(a+b=c\). It remains open.
Stewart and Yu obtained \(H\ll R^{1/3}(\log R)^3\)
\cite[Theorem~1]{SY2001}. Pasten and Sep\'ulveda-Manzo obtained
subexponential estimates under size conditions on the smaller summand
\cite[Theorems~1.2 and~1.3]{PastenSep}. Zhang gives
\(H\ll R^{1/3}\exp(-\log R/(3000\log\log R))\) for all sufficiently large
\(R\) \cite[Theorem~1]{ZhangThird} and subsequently obtains
\(H\ll_\varepsilon R^{2/9+\varepsilon}\) \cite[Theorem~1.1]{Z29}.

\begin{theorem}\label{thm:main}
Let \(\varepsilon>0\). There is a constant \(C_\varepsilon>0\) such that
\[
       \log c\le C_\varepsilon\rad(abc)^\varepsilon
\]
for all coprime positive integers \(a,b\) and \(c=a+b\).
\end{theorem}

Fix \(e=\min(\varepsilon,1/10)\) and suppose \(H\ge R^e\).
For large \(R\), the side of least radical contains a prime power of logarithm \(>2eH\).
The ratio of the other sides is then close to a sign at that prime.
Adjoining auxiliary primes gives a lattice of principal units with large first minimum.
Two determinant estimates force a nonzero polynomial to vanish on a slice of an integer simplex.

Write \(\lambda\) for the slope of the slice.
The first minimum excludes small denominators for its Fourier frequencies.
Steuding and Technau's prime estimate locates a prime \(q\) with
\(q\lambda\) close to \((0,1/m,\ldots,(m-1)/m)\)
\cite[Theorem~3]{SteudingTechnau}.
In each residue class modulo \(q\), the slice thus contains consecutive colors of a smaller simplex.
Finite differences give interpolation on those colors.
Lucas's congruence lifts interpolation to the full slice, contradicting the vanishing polynomial.

The use of AI tools is described at the end of the paper.

